% ============================================================
% Capítulo: Introducción (Bloque B.4)
% ============================================================
\chapter{Introducción}
\label{cap:introduccion}

\section{Contexto y motivación}

La gestión de una biblioteca institucional o municipal -- control de
catálogo, préstamos, devoluciones, reservaciones y multas -- se sigue
resolviendo con frecuencia mediante procesos manuales o herramientas
genéricas no pensadas para el dominio bibliotecario (hojas de cálculo,
registros en papel, o sistemas de escritorio aislados sin acceso
remoto). En el contexto concreto de una biblioteca universitaria como la
de la Universidad Técnica Estatal de Quevedo (UTEQ), esto se traduce en
fricción operativa real y verificable en el propio dominio del
proyecto: un bibliotecario que registra un préstamo a mano no tiene
forma automática de saber si un usuario ya tiene multas pendientes ni de
decrementar el stock disponible de forma atómica; un lector no tiene
forma de autoservicio para ver sus propios préstamos, reservar un libro
o consultar su historial sin acudir físicamente al mostrador.

\textbf{Nota de honestidad (estimación razonada, no dato de fuente
externa verificada).} Este documento evita deliberadamente citar una
cifra de "población afectada" (número de usuarios de biblioteca,
volumen de préstamos anuales, etc.) porque el equipo no cuenta con un
estudio o fuente institucional verificada que la respalde -- inventar
una cifra cuantitativa aquí sería exactamente el tipo de dato fabricado
que el resto de este proyecto se niega a producir (ver, por ejemplo, el
criterio ya aplicado en \texttt{docs/checklists/incose2023-req.md} y en
las notas de honestidad de \texttt{docs/requisitos/SRS-v1.0.0.md}). Lo
que sí se sostiene con evidencia real es el perfil de carga: una
biblioteca universitaria como la de UTEQ tiene un volumen de operación
bajo y no sostenido en el tiempo (ráfagas puntuales en fechas concretas
del calendario académico, no tráfico constante de alta concurrencia),
supuesto ya declarado explícitamente en
\texttt{docs/arquitectura/ISO25010.md} y usado para justificar
decisiones de arquitectura (p.~ej. la prioridad \emph{Media}, no
\emph{Alta}, de la característica "Eficiencia de desempeño").

SGB-SaaS se construye como respuesta a ese problema concreto: una
plataforma web (no de escritorio) que digitaliza el ciclo completo de
préstamo/devolución/reserva/multa, con control de acceso basado en
roles (RBAC) para los distintos perfiles reales de una biblioteca
(lector, bibliotecario, gerente, administrador), y que -- a diferencia
de una implementación mínima que solo "funcione" -- se desarrolla junto
con un programa explícito de evaluación de calidad: pruebas de carga
reales, auditoría de seguridad OWASP Top~10 en vivo, ingeniería de
requisitos formal con trazabilidad verificable, y un protocolo de
usabilidad (SUS) con participantes reales.

\section{Preguntas de investigación}
\label{sec:rqs}

% BORRADOR -- estas 4 preguntas se proponen a partir de los 4 frentes de
% evaluación empírica que el sistema YA ejecuta o tiene protocolo listo
% para ejecutar (rendimiento bajo cache, seguridad OWASP, usabilidad
% SUS, arquitectura híbrida CRUD/SP). Quedan sujetas a validación
% conjunta del equipo antes de considerarse definitivas -- no se tratan
% todavía como un hallazgo cerrado de este documento.

\begin{itemize}[leftmargin=1.2cm]
  \item \textbf{RQ1.} ¿En qué medida el uso de caché (Redis) reduce la
    latencia observable del endpoint de catálogo bajo carga concurrente
    controlada, y con qué margen se cumplen los umbrales de rendimiento
    exigidos?
  \item \textbf{RQ2.} ¿Qué proporción de los controles auditados del
    OWASP Top~10:2021 presenta hallazgos reales en la implementación, y
    en qué proporción de esos hallazgos el propio equipo aplicó una
    corrección verificable dentro del mismo ciclo de desarrollo?
  \item \textbf{RQ3.} ¿Cómo perciben usuarios reales (con roles
    representativos del dominio, p.~ej. lector y bibliotecario) la
    usabilidad del sistema, medida con el instrumento estandarizado
    System Usability Scale (SUS)?
  \item \textbf{RQ4.} ¿Qué efecto tiene adoptar una estrategia híbrida
    de acceso a datos (CRUD elemental vía ORM, operaciones multi-tabla
    vía procedimientos almacenados) sobre la trazabilidad de requisitos
    y la cobertura de pruebas automatizadas, frente a un enfoque de solo
    ORM?
\end{itemize}

\section{Objetivos}

\subsection{Objetivo general}

Evaluar empíricamente la calidad del sistema SGB-SaaS -- en sus
dimensiones de rendimiento, seguridad, usabilidad y arquitectura de
acceso a datos -- mediante evidencia reproducible, versionada y
verificable de forma independiente en el propio repositorio del
proyecto.

\subsection{Objetivos específicos}

\begin{enumerate}[leftmargin=1.2cm]
  \item \textbf{OE1} (trazado a RQ1). Medir, mediante pruebas de carga
    controladas (k6, 50~VUs concurrentes, múltiples corridas
    independientes), la latencia p95 del endpoint de catálogo con y sin
    caché activo, y contrastar el resultado agregado contra los
    umbrales de aceptación exigidos.
  \item \textbf{OE2} (trazado a RQ2). Auditar en vivo, contra el stack
    Docker real (no simulado), un subconjunto representativo de
    controles del OWASP Top~10:2021, documentando cada hallazgo con
    evidencia cruda y aplicando corrección verificable a los defectos
    reales detectados dentro del propio ciclo del proyecto.
  \item \textbf{OE3} (trazado a RQ3). Aplicar el cuestionario System
    Usability Scale (SUS) a una muestra de participantes reales, bajo un
    protocolo de consentimiento informado ya definido, para obtener una
    puntuación de usabilidad percibida con validez estadística.
  \item \textbf{OE4} (trazado a RQ4). Documentar y trazar, mediante
    Architecture Decision Records (ADR) y una matriz de trazabilidad
    verificada automáticamente en CI, la estrategia híbrida de acceso a
    datos, cuantificando qué proporción de los requisitos del sistema
    depende de cada mecanismo (ORM vs.\ procedimiento almacenado) y qué
    porcentaje cuenta con prueba automatizada real.
  \item \textbf{OE5} (transversal a RQ1--RQ4). Construir y versionar
    toda la evidencia empírica del proyecto -- scripts de medición,
    datos crudos, análisis estadístico -- de forma reproducible, para
    que un tercero pueda regenerarla de forma independiente sin
    depender de capturas de pantalla ni de afirmaciones sin respaldo.
\end{enumerate}

\section{Contribuciones}

Este trabajo entrega, con evidencia real ya versionada en el
repositorio al momento de escribir este capítulo:

\begin{itemize}[leftmargin=1.2cm]
  \item Un sistema web funcional y completo (tag \texttt{v1.0.0}) con
    43 requisitos especificados y trazados
    (\texttt{docs/trazabilidad/matriz.csv}: 28 funcionales, 15 no
    funcionales), de los cuales el 100\,\% de los requisitos de
    prioridad \texttt{Must} cuenta con evidencia de verificación real
    (ver \autoref{cap:ingenieria-requisitos}).
  \item Evidencia empírica reproducible y versionada como código, no
    como capturas de pantalla editadas a mano: scripts de prueba de
    carga (\texttt{k6/}), análisis estadístico con test pareado de
    Wilcoxon y tamaño de efecto de Cliff's delta
    (\texttt{scripts/perf-analysis.py}), cobertura de código real
    (JaCoCo) y auditoría de accesibilidad (Lighthouse).
  \item Un artefacto de software citable con identificador persistente:
    el software está archivado en Zenodo con DOI
    (\texttt{10.5281/zenodo.21712467} para la versión \texttt{v0.9.0-rc};
    la versión \texttt{v1.0.0} de este documento se re-archivará al
    cierre de esta entrega, ver el placeholder de DOI en la portada).
  \item Un checklist de calidad de requisitos según el marco INCOSE~v4
    (características C1--C15) aplicado con honestidad metodológica
    completa: documenta explícitamente qué requisitos no cumplen cada
    característica y por qué, en vez de reportar cumplimiento universal
    (\texttt{docs/checklists/incose2023-req.md}).
  \item Documentación arquitectónica versionada como código fuente, no
    como imágenes estáticas: el modelo C4 completo (niveles 1--3) en
    Structurizr DSL (\texttt{docs/arquitectura/workspace.dsl}) y 13
    Architecture Decision Records (\texttt{docs/adr/}).
\end{itemize}

\section{Organización del documento}

El resto de este documento se organiza como sigue.
\autoref{cap:marco-teorico} sienta las bases conceptuales estrictamente
necesarias para el resto del documento -- ingeniería de requisitos según
ISO/IEC/IEEE~29148:2018, el modelo C4, ISO/IEC~25010, REST, autenticación
JWT, OWASP Top~10 y los patrones de acceso a datos (ORM, procedimientos
almacenados, \emph{cache-aside}) --, remitiendo a cada capítulo posterior
para su aplicación concreta. \autoref{cap:trabajos-relacionados} sitúa a
SGB-SaaS frente a la literatura académica indexada disponible sobre
sistemas de gestión bibliotecaria y los patrones arquitectónicos que el
sistema efectivamente utiliza. \autoref{cap:ingenieria-requisitos} resume
el proceso completo de ingeniería de requisitos -- stakeholders, técnicas
de elicitación, los 43 requisitos categorizados, su validación empírica y
su evaluación según el marco de calidad INCOSE~v4.
\autoref{cap:materiales-metodos} describe la metodología de \emph{Design
Science Research} adoptada, los instrumentos de medición usados
(SUS, k6, Lighthouse, JaCoCo, auditoría OWASP), el enfoque de métricas
Goal-Question-Metric y el protocolo experimental replicable del bloque de
rendimiento. \autoref{cap:diseno-arquitectura} describe el diseño
arquitectónico del sistema mediante el modelo C4 en sus tres niveles, las
decisiones arquitectónicas formalizadas como ADR, y el mapeo contra las
características de calidad de producto de ISO/IEC~25010.
\autoref{cap:implementacion} detalla decisiones críticas de
implementación con evidencia de código real: la estrategia de
configuración paramétrica, el manejo uniforme de errores según RFC~7807,
la estrategia de caché y el transporte seguro del token de sesión.
\autoref{cap:amenazas-validez} analiza críticamente las amenazas a la
validez de constructo, interna, externa y de conclusión de la evidencia
empírica recolectada. Los capítulos de Resultados, Discusión,
Conclusiones y Declaraciones/Anexos se incorporarán en tareas de
redacción posteriores, una vez completada la evidencia que dependen de
ellos (ver la nota de estado en \texttt{docs/informe-final.tex}).
